\documentclass[10pt,a4paper]{amsart}
\usepackage[margin=19mm]{geometry}
\usepackage{amsmath,amssymb,mathtools}
\usepackage[scr=rsfs,cal=euler]{mathalfa}
\usepackage{xfrac}
\usepackage[unicode,hidelinks,bookmarks=false]{hyperref}
\newcommand{\ie}{i.e.\ }
\newcommand{\cf}{cf.\ }
\newcommand{\resp}{respectively\ }
\newcommand{\Subsection}{\subsection}
\providecommand{\nobreakdash}{\nobreak\hbox{-}}
\setlength{\emergencystretch}{2em}
\multlinegap0pt
\usepackage{bm}
\usepackage{bbm}
\usepackage{dsfont}
\usepackage{xspace}
\usepackage{cancel}
\usepackage{mathtools}
\usepackage{cleveref}
\usepackage[shortlabels]{enumitem}
\usepackage{amsthm}
\makeatletter
\@ifundefined{lemma}{\newtheorem{lemma}{Lemma}}{}
\@ifundefined{propenum}{\newtheorem{propenum}{Propenum}}{}
\@ifundefined{proposition}{\newtheorem{proposition}[lemma]{Proposition}}{}
\makeatother
\newcommand{\BF}{{\boldsymbol{F}}}
\newcommand{\Bf}{{\boldsymbol{f}}}
\newcommand{\R}{\mathbb{R}}
\newcommand{\bB}{{\boldsymbol{B}}}
\newcommand{\bU}{{\boldsymbol{U}}}
\newcommand{\bb}{{\boldsymbol{b}}}
\newcommand{\bu}{{\boldsymbol{u}}}
\newcommand{\bv}{{\boldsymbol{v}}}
\newcommand{\bw}{{\boldsymbol{w}}}
\title[Asymptotic stability of Landau solutions to the MHD system a]{Asymptotic stability of Landau solutions to the MHD system and energy decay}
\author{Nicola De Nitti, Yun Wang, Shaoheng Zhang}
\date{Source: arXiv:2604.15910v1 (CC BY 4.0)}
\begin{document}
\maketitle

\noindent\emph{Source and licence.} Asymptotic stability of Landau solutions to the MHD system and energy decay. Version arXiv:2604.15910v1 (CC BY 4.0), distributed under the Creative Commons Attribution 4.0 International licence, as recorded in the arXiv licence metadata for this exact version and shown on the abstract page https://arxiv.org/abs/2604.15910v1. Full rights notice: licenses/arxiv-2604.15910-CC-BY-4.0.txt.\par\smallskip

\noindent\textit{Editorial scope.} Counted unit: the Proposition whose source label is \texttt{lemma:a} in the article (source0.tex lines 785--798, source label \texttt{lemma:a}), delivered together with the article's own complete original proof of it (source0.tex lines 800--840, one \texttt{proof} environment of 41 lines). No mathematics is written, repaired or re-derived: statement and proof are the article's own text line for line. Required context retained uncounted: eqn:linearoperator (lines 581--586); prop:decay (lines 711--726); it:lmdecay1 (lines 713--725); eqn:decay1 (lines 715--717); it:lmdecay2 (lines 715--717). Every definition, display and label of those lines is retained unchanged. Omitted from this package and not redistributed: 1--580, 587--710, 718--784, 799--799, 841--1191. Nothing inside a retained excerpt is omitted or altered: every retained body is delivered exactly as the article's lines are written, so no omission receipt of any kind accompanies this package. Citations: the retained excerpts carry no citation command at all, so no bibliography entry of the article is transcribed and no reference is redistributed. Reference adaptation: none was needed and none was made; every reference inside the delivered unit resolves inside it, so no unresolved reference or citation is produced. Printed slips kept literal: no printed word, symbol, hypothesis or number of the counted statement or of its proof was altered. Layout and class: the article's own class is \texttt{amsart} with packages including fontenc, inputenc, babel, csquotes, geometry, graphicx, biblatex, caption, hyperref, cleveref, amsmath, amsthm, amssymb, esint; the wrapper is the lane's pinned \texttt{amsart} editorial wrapper at 10pt on A4, which restates the theorem-like environments the retained text needs and adds the article's own macro definitions that the retained text needs (\texttt{\textbackslash BF}, \texttt{\textbackslash Bf}, \texttt{\textbackslash R}, \texttt{\textbackslash bB}, \texttt{\textbackslash bU}, \texttt{\textbackslash bb}, \texttt{\textbackslash bu}, \texttt{\textbackslash bv}, \texttt{\textbackslash bw}), with extra wrapper packages (bm, bbm, dsfont, xspace, cancel, mathtools, cleveref, enumitem). The delivered numbering is the unit's own. Assets: no retained span contains a figure environment, a graphics-inclusion command, a PDF-page inclusion, a table or a TikZ picture; no image file is copied into this package, the package declares no asset dependency and no image file is redistributed with it. Licence: version arXiv:2604.15910v1 of this article is distributed under the Creative Commons Attribution 4.0 International licence, as recorded in the arXiv licence metadata for this exact version and shown on the abstract page https://arxiv.org/abs/2604.15910v1. A full rights notice is delivered at licenses/arxiv-2604.15910-CC-BY-4.0.txt. Characters: every retained excerpt is pure ASCII, so the delivered engine, XeLaTeX with its default Latin Modern text font, reports no missing character.
\begin{align}\label{eqn:linearoperator}
\begin{aligned}
L_1 \bw&\coloneqq \mathbb P \bigl[- \Delta\bw +  \bU^\bb\cdot \nabla\bw + \nabla\cdot  (\bw \otimes \bU^\bb)\bigr],\\
L_2 \bB&\coloneqq \mathbb P \bigl[- \Delta \bB + \bU^\bb \cdot \nabla \bB-\nabla\cdot (\bB \otimes \bU^\bb )\bigr].
\end{aligned}
\end{align}

\begin{proposition}\label{prop:decay} 
Let $L_i$, with $i =1,\,2$, be the linear operators defined in \cref{eqn:linearoperator}. Then the following propositions hold.
\begin{propenum}
\item \label{it:lmdecay1} For any $1 \le q \le 2$, there exists a constant $\delta_1=\delta_1(q)>0$ (independent of $i$) such that if $0<|\bb|\le \delta_1$, then the semigroup $e^{-tL_i}$ satisfies
\begin{align}\label{eqn:decay1}
\|e^{-tL_i} \Bf \|_2 \le Ct^{-\frac{3}{2}(\frac{1}{q}-\frac{1}{2})} \|\Bf\|_q,\quad \text{for all } t>0, \ \Bf \in  L^q_{\sigma}(\mathbb R^3). 
\end{align}

\item \label{it:lmdecay2}
For any $\frac{6}{5} < q \le 2$, there exists $\delta_2=\delta_2(q)>0$ (independent of $i$) such that if $0<|\bb|\le \delta_2$, then the semigroup $e^{-tL_i}$ satisfies
\begin{align}\label{eqn:decay2}
\|e^{-tL_i}\mathbb P \nabla\cdot \BF\|_2 
\le Ct^{ -\frac{3}{2}(\frac{1}{q}-\frac{1}{2}) -\frac{1}{2} } \| \BF \|_q,\quad \text{for all } t>0, \  \BF\in L^q(\mathbb R^3)^{3\times 3}. 
\end{align}
\end{propenum}
\end{proposition}

\begin{proposition} \label{lemma:a}
Let $L_i$, with $i =1,\,2$, be the linear operators defined in \cref{eqn:linearoperator}. Then the following propositions hold.
\begin{propenum}
\item \label{it:lma1} There exists a constant $\delta_3>0$ such that if $0<|\bb|\le \delta_3$, then for every $\Bf\in L^2_{\sigma}$,
\begin{align}\label{eqn:linearpart}
    \lim\limits_{t\to\infty} \frac{2}{t} \int_{\frac{t}{2}}^t \|e^{-sL_i}\Bf\|_2 \, \mathrm d s=0.
\end{align}    
 \item \label{it:lma2} Let $\frac{6}{5}<r\le2$ and let $\delta_2 = \delta_2(r) > 0$ be the constant from \cref{it:lmdecay2}. If $0<|\bb|\le \delta_2$, then for all $\bu,\bv\in H^1(\R^3)$,
    \begin{align}\label{eqn:nonlinearpart}
    \|e^{-sL_i} \mathbb{P} \nabla\cdot (\bu\otimes \bv)\|_2 \le C s^{-\frac{3}{2}\left(\frac{1}{r}-\frac{1}{2}\right)-\frac{1}{2}} 
    (\|\bu\|_2\|\bv\|_2)^{\frac{3}{2r}-\frac{1}{2}} (\|\nabla\bu\|_2\|\nabla \bv\|_2)^{\frac{3}{2}(1-\frac{1}{r})}.
    \end{align}
\end{propenum}
\end{proposition}

\begin{proof} The two claims are consequences of \cref{prop:decay}.
\begin{enumerate}[label={(\roman*)}]
    \item Let $\delta_3=\min\{\delta_1(\frac{3}{2}),\delta_1(2)\}$, where $\delta_1(r)$ is the constant defined in \cref{it:lmdecay1}. 
For all $\Bf \in L^2_\sigma$ and every $\varepsilon > 0$, we can choose $\Bf_\varepsilon \in C_{c,\sigma}^{\infty}$ such that $\|\Bf - \Bf_\varepsilon\|_2 \le \varepsilon$. Then 
\[
\|e^{-sL_i} \Bf\|_2 
\le \|e^{-sL_i} (\Bf - \Bf_\varepsilon)\|_2 + \|e^{-sL_i} \Bf_\varepsilon\|_2 
\le C \varepsilon + C s^{-\frac{1}{4}} \|\Bf_\varepsilon\|_{\frac32},
\]
where in the last inequality we used \cref{eqn:decay1}. Therefore we have
\[
\frac{2}{t} \int_{\frac{t}{2}}^{t} \|e^{-sL_i} \Bf\|_2 \, \mathrm d s 
\le C\varepsilon + Ct^{-\frac{1}{4}} \|\Bf_\varepsilon\|_{\frac{3}{2}}.
\]
From this, we deduce that
\[
\lim_{t \to \infty} \frac{2}{t} \int_{\frac{t}{2}}^{t} \|e^{-sL_i} \Bf\|_2 \, \mathrm ds \le C \varepsilon.
\]
Since $\varepsilon$ is arbitrary, it follows that $\lim\limits_{t \to \infty} \frac{2}{t} \int_{\frac{t}{2}}^{t} \|e^{-sL_i} \Bf\|_2 \, \mathrm ds = 0$.\\

\item  Using the decay property of $e^{-tL_i}$ in \cref{it:lmdecay2}, we have the estimate
\[
\|e^{-sL_i} \mathbb{P} \nabla\cdot (\bu \otimes \bv)\|_2 
\leq C s^{-\frac{3}{2}\left(\frac{1}{r}-\frac{1}{2}\right) - \frac{1}{2}} \|\bu \otimes \bv\|_r.
\]
H\"older's inequality and Gagliardo--Nirenberg's inequality imply 
\[
\|\bu \otimes \bv\|_r 
\le \|\bu\|_{2r} \| \bv\|_{2r}
\le (\|\bu\|_{2}\|\bv\|_{2})^{1-\theta} 
(\|\nabla \bu\|_{2}\|\nabla \bv\|_{2})^{\theta},
\]
where $\theta$ satisfies $\frac{1}{2r}=\frac{1-\theta}{2}+\theta(\frac{1}{2}-\frac{1}{3})$; in particular, $\theta\in(0,1)$ precisely when $r\in(\frac{6}{5},2]$. Consequently,
\[
\|\bu \otimes \bv\|_r 
\le  (\|\bu\|_{2}\|\bv\|_{2})^{\frac{3}{2r}-\frac{1}{2}} 
(\|\nabla \bu\|_{2}\|\nabla \bv\|_{2})^{\frac{3}{2}(1-\frac{1}{r})},
\]
which completes the proof.
\end{enumerate}
\end{proof}

\end{document}
